\section{Insights} \label{Sec:Insights}

\paragraph{Insight 1: Global-Contract Semantic Failure.}
As shown in Section~\ref{sec:category_correctness}, task category explains nearly three times more variance in semantic correctness than method identity (9.4\% vs.\ 3.3\%), and the correctness gap between easy and hard categories persists despite non-trivial compile rates---ruling out front-end syntax as the primary cause. Static complexity proxies correlate only weakly with failure ($r \leq 0.21$), confirming that the boundary is not reducible to code length or branching depth.
Instead, failures concentrate on tasks where correctness depends on maintaining consistent tensor semantics across different dimensions, memory layouts, and parallel program instances. These constraints span the entire kernel and are therefore difficult to recover through local edits alone.
Case~\ref{sec:case_global_fail} concretizes this mechanism: models produce individually correct Triton idioms while violating the global contract those idioms must collectively satisfy. 
Case~\ref{sec:case_success} provides the contrasting bound: when data dependence is lane-local, no such global contract exists, and all methods are reliably correct.


\paragraph{Insight 2: Repair-Biased Iterative Refinement.}
\label{sec:insight2}

As shown in Section~\ref{sec:iterative_refinement}, iterative refinement reliably expands compilability and correctness, but kernel performance often fails to improve and can even degrade across iterations.
The edit distribution over 352 GEAK diffs confirms that current loops operate primarily as repair mechanisms, with performance-oriented rewrites being rare.
The underlying asymmetry is structural: repair responds to explicit, local error signals (compilation errors, shape mismatches, failing outputs), while performance improvement requires plan-level decisions about tiling, memory layout, and kernel boundaries, which are not recoverable from the feedback available in current iterative pipelines. Case~\ref{sec:case_iter} illustrates this behavior in a concrete setting. As a result, iterative refinement reliably converges to semantically correct implementations, but not to efficient ones. The feedback signal that drove convergence provided insufficient information about this cost, and the iterative process offers no effective means to address it.

\paragraph{Insight 3: Performance as an Unsolved Frontier.}
Among semantically correct kernels, 46.6\% remain slower than eager PyTorch, and the pooled median speedup is only $1.0008\times$ (Section~\ref{sec:overall_results}). 
Cross-machine compounds this: the max/min speedup ratio reaches $21.4\times$ in the worst case, indicating that correct kernels are often hardware-specific rather than generally efficient.
Taken together with Insights~1 and~2, this suggests that correctness and performance represent distinct frontiers: current methods have partially crossed the correctness boundary, but closing the performance gap will likely require qualitatively different mechanisms, such as explicit hardware-aware search or performance-signal feedback, rather than refinements to existing correctness-driven pipelines.
